% TODO make hw stylesheet
\documentclass[10pt]{article}
\usepackage[letterpaper,top=1in]{geometry}
\usepackage[dvipsnames,svgnames]{xcolor}
\usepackage[shortlabels]{enumitem}
\usepackage[framemethod=TikZ]{mdframed}
\usepackage{amsmath,amssymb,amsthm}
\usepackage[colorlinks]{hyperref}
\usepackage{multicol}
\usepackage{tabularx}

\hypersetup{
	colorlinks=true,
	linkcolor=blue,
	filecolor=magenta,
	urlcolor=magenta,
	pdftitle={MAT 528 HW}
}

\usepackage{mathtools}
\usepackage{thmtools}
\usepackage[textsize=footnotesize]{todonotes}
\usepackage{titling}

\reversemarginpar
\mdfdefinestyle{mdbluebox}{roundcorner=10pt,innerbottommargin=9pt,
	linecolor=blue,backgroundcolor=TealBlue!5,}
\declaretheoremstyle[headfont=\sffamily\bfseries\color{MidnightBlue},
mdframed={style=mdbluebox},]{thmbluebox}

\mdfdefinestyle{mdbloobox}{frametitlefont=\bfseries,innerbottommargin=8pt,
	nobreak=true,backgroundcolor=TealBlue!5,linecolor=blue,}
\declaretheoremstyle[headfont=\bfseries\color{MidnightBlue},
mdframed={style=mdbloobox},headpunct={\\[3pt]},postheadspace=0pt,]{thmbloobox}

\mdfdefinestyle{mdredbox}{frametitlefont=\bfseries,innerbottommargin=8pt,
	nobreak=true,backgroundcolor=Salmon!5,linecolor=RawSienna,}
\declaretheoremstyle[headfont=\bfseries\color{RawSienna},
mdframed={style=mdredbox},headpunct={\\[3pt]},postheadspace=0pt,]{thmredbox}

\mdfdefinestyle{mdgreenbox}{linecolor=ForestGreen,backgroundcolor=ForestGreen!5,
	linewidth=2pt,rightline=false,leftline=true,topline=false,bottomline=false,}
\declaretheoremstyle[headfont=\bfseries\sffamily\color{ForestGreen!70!black},
mdframed={style=mdgreenbox},headpunct={ --- },]{thmgreenbox}

\mdfdefinestyle{mdorangebox}{linecolor=RedOrange,backgroundcolor=RedOrange!5,
	linewidth=2pt,rightline=false,leftline=true,topline=false,bottomline=false,}
\declaretheoremstyle[headfont=\bfseries\sffamily\color{RedOrange!70!black},
mdframed={style=mdorangebox},headpunct={ --- },]{thmorangebox}

\mdfdefinestyle{mdblackbox}{linecolor=black,backgroundcolor=RedViolet!5!gray!5,
	linewidth=3pt,rightline=false,leftline=true,topline=false,bottomline=false,}
\declaretheoremstyle[mdframed={style=mdblackbox}]{thmblackbox}

\declaretheoremstyle[spaceabove=6pt,spacebelow=6pt]{basehead}

\declaretheorem[style=basehead,name=Problem]{problem}
\declaretheorem[style=thmredbox,name=Problem,numbered=no]{problem*}
\declaretheorem[style=thmbloobox,name=Setup]{setup} % for config geo
\declaretheorem[style=thmredbox,name=Example,sibling=problem]{example}
\declaretheorem[style=thmredbox,name=$\bigstar$ Problem,sibling=problem]{niceprob}
\declaretheorem[style=thmbluebox,name=Theorem,sibling=problem]{theorem}
\declaretheorem[style=thmbluebox,name=Corollary,sibling=theorem]{corollary}
\declaretheorem[style=thmbluebox,name=Proposition,sibling=theorem]{prop}
\declaretheorem[style=thmbluebox,name=Lemma,numberwithin=problem]{lemma}
\declaretheorem[style=thmbluebox,name=Theorem,numbered=no]{theorem*}
\declaretheorem[style=thmbluebox,name=Lemma,numbered=no]{lemma*}
\declaretheorem[style=thmgreenbox,name=Claim,numberwithin=problem]{claim}
\declaretheorem[style=thmgreenbox,name=Claim,numbered=no]{claim*}
\declaretheorem[style=thmblackbox,name=Remark,numberwithin=problem]{remark}
\declaretheorem[style=thmblackbox,name=Remark,numbered=no]{remark*}
\declaretheorem[style=thmgreenbox,name=Definition,sibling=theorem]{definition}
\declaretheorem[style=thmgreenbox,name=Definition,numbered=no]{definition*}

\providecommand{\ol}{\overline}
\providecommand{\ceil}[1]{\left\lceil #1 \right\rceil}
\providecommand{\floor}[1]{\left\lfloor #1 \right\rfloor}
\newcommand{\dd}{\mathrm{d}}
\providecommand{\ul}{\underline}
\providecommand{\eps}{\varepsilon}
\providecommand{\half}{\frac{1}{2}}
\providecommand{\CC}{\mathbb C}
\providecommand{\FF}{\mathbb F}
\providecommand{\NN}{\mathbb N}
\providecommand{\QQ}{\mathbb Q}
\providecommand{\RR}{\mathbb R}
\providecommand{\ZZ}{\mathbb Z}
\providecommand{\dg}{^\circ}
\providecommand{\ii}{\item}
\providecommand{\setdiff}{\smallsetminus}
\providecommand{\alert}[1]{{\sffamily\textbf{\textcolor{blue}{#1}}}}
\providecommand{\inv}{^{-1}}
\providecommand{\mb}[1]{\mathbf{#1}}

\newenvironment{soln}{\begin{proof}[Solution]}{\end{proof}}

\providecommand{\pow}{\operatorname{Pow}}
\providecommand{\vocab}[1]{{\textbf{\textcolor{ForestGreen}{#1}}}}
\providecommand{\lcm}{\operatorname{lcm}}
\providecommand{\abs}[1]{\left|#1\right|}

\usepackage{mathrsfs}

% place title at top right
\setlength{\droptitle}{-8ex}
\pretitle{\begin{flushright}\Large\bfseries}
\posttitle{\par\end{flushright}}
\preauthor{\begin{flushright}}
\postauthor{\end{flushright}}
\predate{\begin{flushright}}
\postdate{\end{flushright}}

\title{MAT 528 HW03}
\author{\large Aditya Pahuja}
\date{\large Due 09/18}
\begin{document}
\maketitle

\begin{problem}
    [Munkres 19.6]
    Let $\mb x_1$, $\mb x_2$, \dots be a sequence of the points
    of the product space $\prod_{\alpha\in I} X_\alpha$. Show that this sequence
    converges to the point $\mb x$ if and only if
    the sequence $\pi_\alpha(\mb x_1)$, $\pi_\alpha(\mb x_2)$, \dots
    converges to $\pi_\alpha(\mb x)$ for each $\alpha$.
    Is this fact true if one uses the box topology
    instead of the product topology?
\end{problem}

\begin{soln}
    Let $X = \prod_{\alpha\in I} X_\alpha$.
    Suppose that $\mb x_1, \mb x_2, \dots\in X$ converges to $\mb x$
    and fix an arbitrary $\alpha\in I$.
    If $V\subseteq X_\alpha$ is an open set containing $\pi_\alpha(\mb x)$
    then $U = \pi_\alpha\inv(V)\subseteq X$ is an open set
    (by definition of the product topology) that contains $\mb x$.
    Then $U$ contains cofinitely many of the $\mb x_j$,
    so $\pi_\alpha(U) = V$ contains cofinitely many of the $\pi_\alpha(\mb x_j)$.
    Since $V$ was an arbitrary open set, this means that any open set
    $V\subseteq X_\alpha$ containing $\pi_\alpha(\mb x)$
    contains cofinitely many points in the sequence
    $(\pi_\alpha(\mb x_j))_{j\in\NN}$, which is what it means for
    this sequence to converge to $\pi_\alpha(\mb x)$.

    Conversely suppose that for some sequence $\mb x_1,\mb x_2,\dots\in X$,
    the sequence $\pi_\alpha(\mb x_1), \pi_\alpha(\mb x_2), \dots\in X_\alpha$
    converges to $\pi_\alpha(\mb x)$ for each $\alpha\in I$.
    Let $V\subseteq X$ be a set of the form
    \begin{equation*}
	\pi_{\alpha_1}\inv(V_1)\cap\cdots\cap\pi_{\alpha_n}\inv(V_n)
    \end{equation*}
    for some indices $\alpha_1$, \dots, $\alpha_n$
    and open sets $V_i\subseteq X_{\alpha_i}$.
    For each $i=1,\dots,n$, $V_i$ contains cofinitely many $\pi_{\alpha_i}(\mb x_j)$
    because $V_i\subseteq X_{\alpha_i}$ is an open set containing $\pi_\alpha(\mb x)$,
    hence the pre-images of each $V_i$ contains cofinitely many $\mb x_j$.
    Equivalently, $\pi_{\alpha_i}(V_i)^c$ contains finitely many
    points $\pi_{\alpha_i}(\mb x_j)$, so their union
    \begin{equation*}
	\pi_{\alpha_1}(V_1)^c\cup\cdots\cup\pi_{\alpha_n}(V_n)^c
    \end{equation*}
    also contains finitely many points of the sequence $(\mb x_j)_{j\in\NN}$.
    Therefore the complement of this union,
    \begin{equation*}
	\pi_{\alpha_1}(V_1)\cap\cdots\cap\pi_{\alpha_n}(V_n),
    \end{equation*}
    contains cofinitely many of the $\mb x_j$.
    Sets of this form (intersections of finitely pre-images of open sets
    in some $X_\alpha$s) are a basis for the product topology on $X$,
    so any open set $U$ with $\mb x$ as an element contains such a set $B$;
    since only finitely many of the $\mb x_j$ are outside $B$,
    only finitely many of the $\mb x_j$ can be outside the bigger set $U$.
    Since $U$ is arbitrary we have shown that any open set $U$ containing $\mb x$
    also contains cofinitely many $\mb x_j$,
    hence $\mb x_1$, $\mb x_2$, \dots converges to $\mb x$.
    \\

    The problem statement does not hold in the box topology.
    Consider the sequence $\mb x_n = (1/n, 1/n, \dots)\in\RR^\omega$
    (where of course $\RR^\omega$ has the box topology
    inherited from giving each copy of $\RR$ the standard topology).
    Then, projecting this sequence into any of the copies of $\RR$
    gives a sequence converging to zero, so if the problem statement were to hold
    one should expect the $\mb x_n$ to converge to $(0,0,\dots)$.
    However, the neighborhood $\prod_{n=1}^\infty (-1/n,1/n)$ of
    $(0,0,\dots)$ does not contain any of the $\mb x_n$ at all,
    so the $\mb x_n$ don't converge to $(0,0,\dots)$.
\end{soln}
\pagebreak

\begin{problem}
    [Munkres 20.3]
    Let $X$ be a metric space with metric $d$.
    \begin{enumerate}[(a)]
	\ii Show that $d\colon X\times X\to\RR$ is continuous.
	\ii Let $X'$ denote a space with the same underlying set as $X$.
	Show that if $d\colon X'\times X'\to\RR$ is continuous,
	then the topology of $X'$ is finer than the topology of $X$.
\end{enumerate}
\end{problem}

\begin{soln}
    (a) Since the topology of $X$ is generated by the basis of open balls
    $B_r(a) = \{x\in X : d(a,x)<r\}$ for $a\in X$, $r\in(0,\infty)$,
    the topology of $X\times X$ has basis
    \begin{equation*}
	\mathcal B = \{ B_{r_1}(a_1)\times B_{r_2}(a_2) :
	r_1,r_2\in(0,\infty), a_1,a_2\in X\}.
    \end{equation*}
    To show that $d$ is continuous, it suffices to verify that
    for every open interval $(p,q)\subseteq\RR$ with $p<q$,
    $d\inv((p,q))$ is open in $X\times X$;
    it is sufficient because the open intervals are a basis for the topology of $\RR$.
    
    To show that $d\inv((p,q))$ is open in $X\times X$,
    we want to show that for any element of $(x,y)\in d\inv((p,q))$
    there is an element $B$ of $\mathcal B$ such that
    $(x,y)\in B\subseteq d\inv((p,q))$.
    Let $r = d(x,y)$ and define $\delta = \frac{\min(r-p, q-r)}{2}$.
    Then I claim that if $(x',y')\in B = B_\delta(x)\times B_\delta(y)$,
    $d(x',y')\in (p,q)$, or equivalently $(x',y')\in d\inv((p,q))$.
    By a few applications of the triangle inequality, we get the bounds
    \begin{align*}
        d(x',y') \le d(x',x) + d(x,y') \le d(x',x) + d(x,y) + d(y,y')
	< r + 2\delta \le r + 2\cdot\frac{q-r}2 = q\\
        d(x',y') \ge d(x,y') - d(x,x') \ge d(x,y) - d(y,y') - d(x,x')
	> r - 2\delta \ge r - 2\cdot\frac{r-p}2 = p.
    \end{align*}
    Therefore any $(x',y')\in B$ is in the pre-image of $(p,q)$,
    so $B\subseteq d\inv((p,q))$. That is, for every $x\in d\inv((p,q))$,
    we have found a basis element containing $x$ and contained in $d\inv((p,q))$,
    which implies that $d\inv((p,q))$ is open in $X\times X$
    for all open intervals $(p,q)$,
    so $d\colon X\times X\to\RR$ is continuous.
    \\

    (b) $d$ is continuous in each coordinate
    so for any fixed $a\in X'$,
    $d_a = d|_{\{a\}\times X}\colon X'\to\RR$ is continuous.
    Then for any $r\in\RR_{>0}$, the pre-image $d_a\inv((-r,r))$ is
    $\{(a, x): d(a,x) < r\}$.

    \begin{lemma}
        \label{lem:embedding-X-in-X2}
	$U\subseteq X'$ is open if and only if $\{a\}\times U\subseteq \{a\}\times X'$
	is open (where $\{a\}\times X'$ has the subspace topology).
    \end{lemma}
    \begin{proof}
        Suppose that $U\subseteq X'$ is open.
	Then $\{a\}\times U = (X'\times U)\cap (\{a\}\times X')$
	is the intersection of an open subset $X'\times U\subseteq X'\times X'$
	(because products of open sets are open in the product topology
	on $X'\times X'$)
	with $\{a\}\times X'$, which means it's open in
	the subspace topology on $\{a\}\times X'$.

	Conversely suppose $\{a\}\times U$ is open in $\{a\}\times X'$.
	Then there is an open set $V\subseteq X'\times X'$
	such that $V\cap(\{a\}\times X') = \{a\}\cap U$.
	Then, $V\cap(X'\times U)$ is also open in $X'$
	and the projection $U' = \pi_2(V\cap(X'\times U))$ in the second coordinate
	is an open set because $\pi_2$ is an open map due to HW1 \#5.
	On one hand $U'$ is contained in $U$ because it is a subset of
	$\pi_2(X'\times U) = U$, and on the other hand it contains $U$
	because $V$ contains every point of the form $(a,u)$ for some $u\in U$,
	so the open set $U'$ is equal to $U$.
    \end{proof}
    By the lemma, the map $p\colon X'\to\{a\}\times X'$ given by $p(x) = (a,x)$
    is a homeomorphism: it is bijective because it has an inverse $p\inv(a,x) = x$,
    it is continuous because the pre-image under $p$ of an open
    $\{a\}\times U\subseteq\{a\}\times X'$ is the open (by the Lemma) set $U$,
    and its inverse is continuous because
    the pre-image of an open $U$ under $p\inv$ is $\{a\}\times U$
    which is again open by the Lemma.

    Now we are set to solve the problem: the set
    \begin{equation*}
	p(d_a\inv((-r,r))) = p(\{(a,x)\in\{a\}\times X' : d(a,x) < r\}) =
	\{x\in X' : d(a,x) < r\} = B_r(a)
    \end{equation*}
    is open in $X'$ by the work above, which means that
    every basis element of $X$ (the open balls of $d$) is open in $X'$,
    so the topology generated by the basis of $X$, i.e. the topology of $X$,
    is coarser than that of $X'$ because the generated topology
    is the minimal one containing the topology of $X$.
\end{soln}

\pagebreak

\begin{problem}
    [Munkres 20.10]
    Let $X$ denote the subset or $\RR^\omega$ consisting of
    all sequences $(x_1,x_2,\dots)$ such that $\sum x_i^2$ converges.
    \begin{enumerate}[(a)]
	\ii Show that if $\mb x,\mb y\in X$, then $\sum\abs{x_iy_i}$ converges.
	\ii Let $c\in\RR$. Show that if $\mb x,\mb y\in X$,
	then so are $\mb x+\mb y$ and $c\mb x$.
	\ii Show that
	\begin{equation*}
	    d(\mb x,\mb y) = \sqrt{\sum_{i=1}^\infty (x_i-y_i)^2}
	\end{equation*}
	is a well-defined metric on $X$.
\end{enumerate}
\end{problem}

\begin{soln}
    (a) We know that $\sum_{i=1}^\infty x_i^2$ and $\sum_{i=1}^\infty y_i^2$
    both converge absolutely because all their addends are nonnegative,
    so the series $\sum_{i=1}^\infty (x_i^2 + y_i^2)$ is also (absolutely) convergent.
    We can then bound $x_i^2 + y_i^2 \ge 2\abs{x_iy_i} \ge \abs{x_iy_i}$,
    so since $\abs{x_iy_i}$ is also nonnegative we can conclude that
    \begin{equation*}
	\sum_{i=1}^\infty \abs{x_iy_i}\le\sum_{i=1}^\infty (x_i^2+y_i^2)<\infty
    \end{equation*}
    is convergent.
    \\

    (b) We need to show that $\sum_{i=1}^\infty (x_i+y_i)^2$ converges.
    This follows from rearranging the sum as
    \begin{equation*}
	\sum_{i=1}^\infty (x_i + y_i)^2 =
	\sum_{i=1}^\infty x_i^2 + \sum_{i=1}^\infty 2x_iy_i + \sum_{i=1}^\infty y_i^2
    \end{equation*}
    where the first and last sum converge by $\mb x,\mb y\in X$
    and the middle sum converges because $\sum_{i=1}^\infty \abs{x_iy_i}$
    converges, and of course doubling all the addends of a convergent series
    yields another convergent series.
    This proves $\mb x+\mb y\in X$.

    Then, given $c\in\RR$ and $\mb x\in X$, we find that
    \begin{equation*}
	\sum_{i=1}^\infty (cx_i)^2 = c^2\sum_{i=1}^\infty x_i^2
    \end{equation*}
    is also convergent so $c\mb x$ is in $X$ as well.
    \\

    (c) Note that
    \begin{equation*}
	\mb x\cdot \mb y = \sum_{i=1}^\infty x_iy_i
    \end{equation*}
    converges if $\mb x,\mb y\in X$ by part (a)
    and is an inner product on $X$, and the norm of $\mb x$
    induced by this inner product is $\sqrt{\mb x\cdot \mb x} = d(\mb x,\mb 0)$.
    By the Cauchy-Schwarz inequality,
    $\abs{\mb x\cdot\mb y}\le d(\mb x,\mb 0)d(\mb y,\mb 0)$, so
    \begin{align*}
	d(\mb x,\mb y)^2
	&= (\mb x-\mb y)\cdot(\mb x-\mb y)\\
	&= d(\mb x,0)^2 + d(\mb y,0)^2 - 2\mb x\cdot \mb y\\
	&\le d(\mb x,0)^2 + d(\mb y,0)^2 + 2d(\mb x,\mb 0)d(\mb y,\mb 0)\\
	&= (d(\mb x,\mb 0) + d(\mb y,\mb 0))^2.
    \end{align*}
    Since the quantities being squared are both nonnegative
    we can safely take the square root to recover the triangle inequality
    when the point in the middle is zero;
    this is sufficient because
    \begin{equation*}
	d(\mb x,\mb z) + d(\mb z,\mb y) =
	d(\mb x-\mb z,\mb0) + d(\mb0,\mb y-\mb z)
	\ge d(\mb x-\mb z,\mb y-\mb z) = d(\mb x,\mb y).\qedhere
    \end{equation*}
\end{soln}

\begin{problem}
    [Munkres 21.2]
    Let $X$ and $Y$ be metric spaces with metrics $d_X$ and $d_Y$, respectively.
    Let $f\colon X\to Y$ have the property that
    for every pair of points $x_1,x_2\in X$,
    \begin{equation*}
	d_X(x_1,x_2) = d_Y(f(x_1), f(x_2)).
    \end{equation*}
    Show that $f$ is an embedding.
    It is called an \vocab{isometric embedding} of $X$ in $Y$.
\end{problem}

\begin{soln}
    $f$ is injective (so $f\colon X\to f(X)$ is bijective) because
    \begin{equation*}
        f(x_1)=f(x_2) \iff d(f(x_1),f(x_2)) = 0 \iff
	d(x_1,x_2)=0 \iff x_1=x_2.
    \end{equation*}
    Moreover $f$ is continuous via the $\eps$-$\delta$ definition of continuity
    because for every $\eps>0$ there is $\delta=\eps>0$ such that
    \begin{equation*}
	d_X(x_1,x_2) < \delta = \eps \iff d_Y(f(x_1),f(x_2)) = d(x_1,x_2) < \eps
    \end{equation*}
    and conversely $f\inv\colon f(X)\to X$ is continuous
    because $y_1,y_2\in f(X)$ we get the relevant $\eps$-$\delta$ bound
    by taking $x_1 = f\inv(y_1)$ and $x_2=f\inv(y_2)$ in the displayed statement.
    Thus $f$ is a homeomorphism between $X$ and $f(X)$,
    which is what it means to be an embedding of $X$ into $Y$.
\end{soln}

\begin{problem}
    [Munkres 21.9]
    Let $f_n\colon\RR\to\RR$ be the function
    \begin{equation*}
	f_n(x) = \frac{1}{n^3(x-1/n)^2 + 1}.
    \end{equation*}
    Let $f\colon\RR\to\RR$ be the zero function.
    \begin{enumerate}[(a)]
	\ii Show that $f_n(x)\to f(x)$ for each $x\in\RR$.
	\ii Show that $f_n$ does not converge uniformly to $f$.
\end{enumerate}
\end{problem}

\begin{soln}
    (a) For a given $x$, the sequence $a_n = n^3(x-1/n)^2 + 1 = n(nx - 1)^2 + 1$
    has limit $\infty$ as $n$ goes to infinity because
    it is a nonconstant polynomial in $n$
    with positive leading coefficient. Therefore the sequence
    of its reciprocals $1/a_n = f_n(x)$ converges to zero.\\

    (b) $f_n$ converges to $f$ uniformly if
    for each $\eps>0$ there is an $n$ such that $n>N$ implies that
    $|f_n(x) - f(x)| = f_n(x) < \eps$ for all $x$.
    However, given $\eps<1$, such an $n$ doesn't exist, as for any $n$,
    $x=1/n$ gives
    \begin{equation*}
	\abs{f_n(1/n) - f(1/n)} = f_n(1/n) = 1 > \eps.
    \end{equation*}
    Therefore, the convergence of $f_n$ to $f$ is not uniform.
\end{soln}










\end{document}

